% Coverage: physical pp. 48--52 (printed pp. 25--29), from the Section 25.1
% heading through the paragraph ending immediately before Section 25.2.
% Equations (25.1.1)--(25.1.20), three unnumbered displays, citation [2], and
% no footnotes, figures, tables, or typographic dividers. Uncertainty: none.

\subsection{Graded Lie Algebras and Graded Parameters}

\providecommand{\gradedbracket}[2]{\mathopen{[}#1,#2\mathclose{\}}}

We saw in Section 2.2 how to express any continuous symmetry transformation in
terms of a Lie algebra of linearly independent symmetry generators $t_a$ that
satisfy commutation relations
$[t_a,t_b]=i\sum_c C^c_{ab}t_c$. In much the same way, supersymmetry is
expressed in terms of symmetry generators $t_a$ that form a \emph{graded} Lie
algebra~\hyperref[ch25-ref-2]{[2]}, embodied in commutation and
anticommutation relations of the form
\begin{equation}
  t_at_b-(-1)^{\eta_a\eta_b}t_bt_a
  =i\sum_c C^c_{ab}t_c\,.
  \tag{25.1.1}\label{eq:25.1.1}
\end{equation}
(The summation convention is suspended in this section.) Here $\eta_a$ for
each $a$ is either $+1$ or 0 and is known as the \emph{grading} of the
generator $t_a$, and the $C^c_{ab}$ are a set of numerical structure
constants. Generators $t_a$ for which $\eta_a=1$ are called \emph{fermionic};
the others, for which $\eta_a=0$, are called \emph{bosonic}.
Eq.~\eqref{eq:25.1.1} provides commutation relations for bosonic operators with
each other and with fermionic operators, but anticommutation relations for
fermionic operators with each other. We will come back soon to the motivation
for Eq.~\eqref{eq:25.1.1}; for the moment we will just take a look at its
consequences for the structure constants.

According to Eq.~\eqref{eq:25.1.1}, the structure constants must satisfy the
conditions
\begin{equation}
  C^c_{ba}=-(-1)^{\eta_a\eta_b}C^c_{ab}\,.
  \tag{25.1.2}\label{eq:25.1.2}
\end{equation}

For any operator formed as a functional of field operators, the products of
two bosonic or two fermionic operators are bosonic, and the products of a
fermionic with a bosonic operator are fermionic, so that
\begin{equation}
  C^c_{ab}=0\quad\text{unless}\quad
  \eta^c=\eta^a+\eta^b\pmod 2\,.
  \tag{25.1.3}\label{eq:25.1.3}
\end{equation}
Also, for any operator formed in this way, the Hermitian adjoint of a bosonic
or fermionic operator is, respectively, bosonic or fermionic. If the $t_a$ are
Hermitian operators, then the structure constants satisfy a reality condition
\begin{equation}
  C^{c*}_{ab}=-C^c_{ba}\,.
  \tag{25.1.4}\label{eq:25.1.4}
\end{equation}

The structure constants also satisfy a non-linear constraint, which follows
from a super-Jacobi identity
\begin{equation}
  (-1)^{\eta_c\eta_a}
    \gradedbracket{\gradedbracket{t_a}{t_b}}{t_c}
  +(-1)^{\eta_a\eta_b}
    \gradedbracket{\gradedbracket{t_b}{t_c}}{t_a}
  +(-1)^{\eta_b\eta_c}
    \gradedbracket{\gradedbracket{t_c}{t_a}}{t_b}
  =0\,.
  \tag{25.1.5}\label{eq:25.1.5}
\end{equation}
Here `$[\ldots\}$' denotes a commutator/anticommutator like that appearing on
the left-hand side of Eq.~\eqref{eq:25.1.1}, but here extended to any graded
operators $O$, $O'$, etc.
\begin{equation}
  \gradedbracket{O}{O'}
  \equiv OO'-(-1)^{\eta(O)\eta(O')}O'O
  =-(-1)^{\eta(O)\eta(O')}\gradedbracket{O'}{O}\,,
  \tag{25.1.6}\label{eq:25.1.6}
\end{equation}
it now being understood that any product $O=t_at_bt_c\cdots$ of generators is
given a grading
$\eta(O)\equiv\eta_a+\eta_b+\eta_c+\cdots\pmod 2$. (To prove
Eq.~\eqref{eq:25.1.5}, it is sufficient to prove that the coefficients of
$t_at_bt_c$ and $t_at_ct_b$ vanish, for then the symmetry of the left-hand
side of Eq.~\eqref{eq:25.1.5} under the cyclic permutations
$abc\longrightarrow bca\longrightarrow cab$ will ensure that the coefficients
of all other products of generators also vanish. The coefficient of
$t_at_bt_c$ in Eq.~\eqref{eq:25.1.5} is
\[
  (-1)^{\eta_c\eta_a}
  -(-1)^{\eta_a\eta_b}(-1)^{\eta_a(\eta_b+\eta_c)}
  =0\,,
\]
while the coefficient of $t_at_ct_b$ is
\[
  (-1)^{\eta_a\eta_b}(-1)^{\eta_b\eta_c}
    (-1)^{\eta_a(\eta_b+\eta_c)}
  -(-1)^{\eta_b\eta_c}(-1)^{\eta_c\eta_a}
  =0\,,
\]
completing the proof.) By inserting Eq.~\eqref{eq:25.1.1} in
Eq.~\eqref{eq:25.1.5}, we find the constraint
\begin{equation}
  \sum_d(-1)^{\eta_c\eta_a}C^d_{ab}C^e_{dc}
  +\sum_d(-1)^{\eta_a\eta_b}C^d_{bc}C^e_{da}
  +\sum_d(-1)^{\eta_b\eta_c}C^d_{ca}C^e_{db}
  =0\,.
  \tag{25.1.7}\label{eq:25.1.7}
\end{equation}
Of course, in the case where all generators are bosonic
Eq.~\eqref{eq:25.1.5} is the usual Jacobi identity, and
Eq.~\eqref{eq:25.1.7} is the usual non-linear condition (2.2.22) on the
structure constants.

Eq.~\eqref{eq:25.1.1} can be taken as our starting point, but it can also be
given a motivation like that given for ordinary Lie algebras in Section 2.2,
as a necessary feature of finite continuous symmetry transformations. The
difference is that now these transformations depend on continuous
\emph{graded} parameters. A set of graded c-number parameters can be thought
of as `numbers,' including Grassmann parameters (see Section 9.5) as well as
ordinary numbers, which satisfy the associative and distributive rules of
arithmetic, but which instead of simply commuting satisfy relations
\begin{equation}
  \alpha^a\beta^b=(-1)^{\eta_a\eta_b}\beta^b\alpha^a\,,
  \tag{25.1.8}\label{eq:25.1.8}
\end{equation}
where $\alpha^a,\beta^a,\ldots$ are used to distinguish different values of
the $a$th parameter, in the same way that in vector algebra we might let
$v^a$ and $u^a$ denote the $a$-components of different values of some real
vector. Again the $a$th graded parameter is given a grading $\eta_a$ equal to
$+1$ or 0 when $\alpha^a$ is a fermionic or bosonic parameter, respectively.
That is, these parameters commute if either is bosonic, and anticommute if
both are fermionic. The product $\alpha^a\beta^b\gamma^c\cdots$ of a set of
graded parameters is given the grading
$\eta_a+\eta_b+\eta_c+\cdots\pmod 2$; that is, such a product is fermionic if
it involves an odd number of fermionic parameters, and is otherwise bosonic.
With this grading, it is easy to see that products of graded parameters
satisfy a commutation or anticommutation rule just like
Eq.~\eqref{eq:25.1.8}.

Consider a continuous transformation $T(\alpha)$, given by a formal power
series in the graded parameters $\alpha^a$:
\begin{equation}
  T(\alpha)=\mathbf{1}+\sum_a\alpha^at_a
  +\sum_{ab}\alpha^a\alpha^bt_{ab}+\cdots\,,
  \tag{25.1.9}\label{eq:25.1.9}
\end{equation}
where $t_a,t_{ab}$, etc.\ are a set of $\alpha$-independent operator
coefficients, not yet assumed to satisfy any algebraic relations like
Eq.~\eqref{eq:25.1.1}. Because the parameters $\alpha^a$ satisfy
Eq.~\eqref{eq:25.1.8}, the coefficients $t_{ab\ldots}$ must satisfy
symmetry/antisymmetry conditions, such as
\begin{equation}
  t_{ab}=(-1)^{\eta_a\eta_b}t_{ba}\,.
  \tag{25.1.10}\label{eq:25.1.10}
\end{equation}
It is convenient also to assume that the transformation $T(\beta)$ commutes
with any value $\alpha^a$ of any graded parameter, in which case the operator
coefficients in \eqref{eq:25.1.9} satisfy the conditions
\begin{equation}
  \alpha^at_b=(-1)^{\eta_a\eta_b}t_b\alpha^a\,,
  \tag{25.1.11}\label{eq:25.1.11}
\end{equation}
\begin{equation}
  \alpha^at_{bc}=(-1)^{\eta_a(\eta_b+\eta_c)}t_{bc}\alpha^a\,.
  \tag{25.1.12}\label{eq:25.1.12}
\end{equation}
That is, $t_b$ and $t_{bc}$ commute or anticommute with graded parameters as
if they were graded parameters themselves, with grading $\eta_b$ and
$\eta_b+\eta_c\pmod 2$, respectively.

The other constraints on these operators follow from the requirement that the
$T(\alpha)$ form a semi-group; that is, that the product of the operators for
different values $\alpha$ and $\beta$ of the graded parameters is itself a
$T$-operator
\begin{equation}
  T(\alpha)T(\beta)=T\bigl(f(\alpha,\beta)\bigr)\,,
  \tag{25.1.13}\label{eq:25.1.13}
\end{equation}
where $f^c(\alpha,\beta)$ is itself a formal power series in the graded
parameters. Because $T(0)T(\beta)=T(\beta)$ and
$T(\alpha)T(0)=T(\alpha)$, we must have
\begin{equation}
  f^c(0,\beta)=\beta^c,\qquad
  f^c(\alpha,0)=\alpha^c\,,
  \tag{25.1.14}\label{eq:25.1.14}
\end{equation}
and therefore the power series expansion for $f(\alpha,\beta)$ must take the
form
\begin{equation}
  f^c(\alpha,\beta)=\alpha^c+\beta^c
  +\sum_{ab}f^c_{ab}\alpha^a\beta^b+\cdots\,,
  \tag{25.1.15}\label{eq:25.1.15}
\end{equation}
where $f^c_{ab}$ is a set of ordinary (that is, bosonic) constants, and
`$\cdots$' denotes terms of third or higher order in the graded parameters. In
order for $f^c(\alpha,\beta)$ to be a graded parameter, it is necessary for
each term in Eq.~\eqref{eq:25.1.15} to have the same grading, which implies
that
\begin{equation}
  f^c_{ab}=0\quad\text{unless}\quad
  \eta^c=\eta^a+\eta^b\pmod 2\,.
  \tag{25.1.16}\label{eq:25.1.16}
\end{equation}
Inserting the power series \eqref{eq:25.1.9} and \eqref{eq:25.1.15} into the
product rule \eqref{eq:25.1.13} gives
\[
  \begin{aligned}
  &\left[\mathbf{1}+\sum_a\alpha^at_a
    +\sum_{ab}\alpha^a\alpha^bt_{ab}+\cdots\right]
   \left[\mathbf{1}+\sum_a\beta^at_a
    +\sum_{ab}\beta^a\beta^bt_{ab}+\cdots\right]
  \\
  &\quad=\mathbf{1}
    +\sum_c\left(\alpha^c+\beta^c
      +\sum_{ab}f^c_{ab}\alpha^a\beta^b+\cdots\right)t_c
  \\
  &\qquad
    +\sum_{cd}\left(\alpha^c+\beta^c+\cdots\right)
      \left(\alpha^d+\beta^d+\cdots\right)t_{cd}+\cdots\,.
  \end{aligned}
\]
The coefficients of $\mathbf{1}$, $\alpha^a$, $\beta^a$,
$\alpha^a\alpha^b$, and $\beta^a\beta^b$ match on both sides of this
equation, but the condition that the coefficients of $\alpha^a\beta^b$ are
the same yields the non-trivial relation
\begin{equation}
  (-1)^{\eta_a\eta_b}t_at_b
  =\sum_c f^c_{ab}t_c+t_{ab}
    +(-1)^{\eta_a\eta_b}t_{ba}
  =\sum_c f^c_{ab}t_c+2t_{ab}\,.
  \tag{25.1.17}\label{eq:25.1.17}
\end{equation}
(The \emph{sign} factor on the left-hand side arises from the interchange of
$t_a$ and $\beta^b$.) Together with higher-order relations of the same sort,
this allows us to calculate the whole function \eqref{eq:25.1.9} if we know
the generators $t_a$ and the group-composition function
$f^a(\alpha,\beta)$. But for this to be possible, $t_a$ must satisfy a
constraint. Using Eq.~\eqref{eq:25.1.10}, the difference or sum of
Eq.~\eqref{eq:25.1.17} and the same equation with $a$ and $b$ interchanged
yields the Lie superalgebra relations Eq.~\eqref{eq:25.1.1}, with structure
constants given by
\begin{equation}
  iC^c_{ab}=(-1)^{\eta_a\eta_b}f^c_{ab}-f^c_{ba}\,.
  \tag{25.1.18}\label{eq:25.1.18}
\end{equation}
Also, Eq.~\eqref{eq:25.1.3} follows immediately from
Eqs.~\eqref{eq:25.1.16} and \eqref{eq:25.1.18}.

The complex conjugate $\alpha^*$ of an anticommuting c-number $\alpha$ is
defined so that the Hermitian adjoint of the product of $\alpha$ and an
arbitrary operator $O$ is
\begin{equation}
  (\alpha O)^\dagger=O^\dagger\alpha^*\,.
  \tag{25.1.19}\label{eq:25.1.19}
\end{equation}
It follows that products of c-numbers behave the same under complex
conjugation as operators do under Hermitian conjugation:
\begin{equation}
  (\alpha\beta)^*=\beta^*\alpha^*\,,
  \tag{25.1.20}\label{eq:25.1.20}
\end{equation}
and that $\alpha^*$ has the same grading as $\alpha$.

The graded Lie algebras of importance to physics are severely restricted by
spacetime symmetries. We now turn to a consideration of these restrictions.
